\chapter{研究现状}

% 本章围绕“如何将人的空间注意可靠地转化为面向特定物体的交互”梳理相关工作。与按通信、视觉和生理信号分别罗列技术不同，本文沿对象选择、身份绑定、意图确认与持续会话这一完整链条组织文献：首先讨论定向光学链路如何把空间方向保留到数据传输中；继而讨论注意力线索如何减少多对象交互中的显式发现与切换；最后讨论眼动与脑电能否进一步区分“看向目标”与“准备操作”。这一组织方式也对应本文从 RetroAR 的“所指即所控”、EchoSight 的“所看即所控”到 Omnisight 的“所思即所控”的递进关系。

\section{定向光学链路}

定向交互首先要求系统回答“当前数据究竟来自哪一个物体”。Wi-Fi、蓝牙等射频技术具有覆盖广和生态成熟的优势，但同一覆盖域中的多个设备共享发现空间，物理对象、网络地址与用户指向之间仍需额外映射。光学无线链路则同时受到发射光束、接收视场和视距关系约束，可以在物理层把通信范围限制到局部方向。因此，光的方向性不仅是信道特征，也能够成为对象选择和身份绑定的交互资源。

早期系统已展示了“以光选物”的基本可行性。Ma et al. 提出的 FindIT Flashlight 使用手持光束唤醒特定标签，使“照向”动作直接参与对象筛选\supercite{findit}；Perli et al. 的 PixNet 通过 LCD--摄像头对构造具有空间分辨能力的链路\supercite{pixnet}。在增强现实场景中，Yang et al. 的 InfoLED 和 Ahuja et al. 的 LightAnchors 进一步复用设备指示灯，将身份信息或虚拟内容锚定到具体光源\supercite{infoled,lightanchors}。这些工作把数字信息附着到空间位置，却主要完成单向发现或显示；对象端主动发光需要额外能量，识别结果也常需再映射到独立的无线控制通道。

\begin{figure}[htbp]
  \centering
  \fbox{\parbox[c][4.0cm][c]{0.86\textwidth}{\centering
  \textbf{图片占位：早期定向光学系统对比}\\[0.6em]
  建议内容：FindIT、PixNet、InfoLED 与 LightAnchors 的收发结构，\\
  对比主动发光、空间分辨、身份承载和后续控制通道。}}
  \caption{早期定向光学链路的代表性系统与功能边界（占位图）}
  \label{fig:early-directional-optical-systems}
\end{figure}

逆反射与光调制为低功耗双向链路提供了另一条路径。Li et al. 提出的 Retro-VLC 将逆反射材料与液晶调制器结合，使标签通过调制入射光回传数据\supercite{retrovlc}；Xu et al. 的 PassiveVLC 通过器件与编码协同设计提高无电池标签的可用性\supercite{passivevlc}；Wu et al. 的 TurboBoosting 则利用多像素液晶和高阶调制提升速率与距离\supercite{turboboost}。这一阶段主要把方向性当作提高空间复用率和降低功耗的信道条件，研究指标集中在误码率、速率、距离和视角，并未系统讨论光束方向与用户所选对象之间的语义对应。

面向密集节点，研究重点又从单链路性能转向空间并发。Xu et al. 的 RetroMUMIMO 挖掘不同液晶标签的脉冲响应差异，在物理层并行解调多个上行并以并发 MAC 降低大规模接入时延\supercite{retromumimo}。2024 年，Ghiasi et al. 同时利用偏振和颜色构造无源可见光 MIMO，使多个液晶子信道的容量接近线性增长\supercite{passivevlc_mimo_2024}；Wang et al. 将事件相机与数字微镜器件结合，并行识别约两千个空间信道，表明事件驱动接收可显著突破传统帧率瓶颈\supercite{selene_2024}。这些进展说明光学空间分集能够提高密集场景下的吞吐与并发度，但“同时读到多个标签”与“只绑定用户此刻关心的标签”仍是两个不同目标。

\begin{figure}[htbp]
  \centering
  \fbox{\parbox[c][4.0cm][c]{0.86\textwidth}{\centering
  \textbf{图片占位：无源光学链路与空间并发}\\[0.6em]
  建议内容：Retro-VLC/PassiveVLC 的单标签链路，及 RetroMUMIMO、\\
  偏振--颜色 MIMO、事件相机+DMD 的多信道设计。}}
  \caption{无源光学链路从单标签通信到空间并发的演进（占位图）}
  \label{fig:passive-optical-concurrency}
\end{figure}

本文的第一项前序工作 RetroAR 正是在这一分界处把通信方向性转化为交互方向性。Li et al. 以商用手机闪光灯和后置摄像头作为读写器，设计动态逆反射标签 ViTag，将手机指向、标签身份读取与六自由度位姿估计放在同一光路中\cite{retroar}。与先做视觉识别再查找 Wi-Fi/BLE 地址的组合方案相比，RetroAR 的关键不在于获得更高通信速率，而在于光束只照亮并读取被指对象，使“物理选择”和“数字身份”同时完成。其系统评估给出最远约 $4\,\mathrm{m}$、最大约 $100^\circ$ 视角的工作范围，用户研究则表明其相较基于 Wi-Fi 的方案显著缩短了多对象任务中的交互时间。

\begin{figure}[htbp]
  \centering
  \fbox{\parbox[c][4.0cm][c]{0.86\textwidth}{\centering
  \textbf{图片占位：无源光学链路与空间并发}\\[0.6em]
  建议内容：Retro-VLC/PassiveVLC 的单标签链路，及 RetroMUMIMO、\\
  偏振--颜色 MIMO、事件相机+DMD 的多信道设计。}}
  \caption{无源光学链路从单标签通信到空间并发的演进（占位图）}
  \label{fig:passive-optical-concurrency}
\end{figure}

第二项前序工作 EchoSight 进一步把移动读写器由手持手机迁移到 AR 眼镜。Li et al. 设计眼镜端 SightReader 与对象端 EchoTag，通过双元件、双波段近红外链路完成指令下行、身份与状态上行，并以对齐检测和按需唤醒抑制非目标响应\supercite{echosight}。该系统不要求预注册或先建立射频连接，达到约 $5\,\mathrm{m}$ 距离和约 $120^\circ$ 视角；其价值在于把“抬起手机并指向”推进为“看向并控制”。不过，EchoSight 的固定光轴实际测量的是眼镜朝向而非眼球视线；当相邻目标密集、头眼分离或视距被遮挡时，目标对齐和会话保持仍会失效。

\begin{figure}[htbp]
  \centering
  \fbox{\parbox[c][4.0cm][c]{0.86\textwidth}{\centering
  \textbf{图片占位：无源光学链路与空间并发}\\[0.6em]
  建议内容：Retro-VLC/PassiveVLC 的单标签链路，及 RetroMUMIMO、\\
  偏振--颜色 MIMO、事件相机+DMD 的多信道设计。}}
  \caption{无源光学链路从单标签通信到空间并发的演进（占位图）}
  \label{fig:passive-optical-concurrency}
\end{figure}

2026 年的研究开始把可见光与射频后向散射置于同一环境物联网架构中，以光完成能量供给、局部发现或控制，再由射频承担非视距连接\supercite{xie_joint_vlc_2026}。这提示定向交互不必把所有阶段都压在单一光路上：光适合确认“是哪一个”，成熟无线链路适合维持“持续可控”。但混合链路也带来新的身份一致性问题，即光学标签与无线地址必须经过可验证绑定，不能把“选中”简单等同于“已授权连接”。

\begin{figure}[htbp]
  \centering
  \fbox{\parbox[c][4.2cm][c]{0.86\textwidth}{\centering
  \textbf{图片占位：定向光学链路的技术演进}\\[0.6em]
  建议内容：主动光标记、无源逆反射、空间并发，\\
  以及 RetroAR 与 EchoSight 从手持指向到佩戴式观察的演进。}}
  \caption{定向光学链路从通信到对象绑定的演进（占位图）}
  \label{fig:directional-optical-links}
\end{figure}

\medskip
\noindent\textbf{本节小结。}现有光学研究已经解决“低功耗地传数据”和“并发地读标签”，RetroAR 与 EchoSight 则进一步把定向光路用于目标绑定，实现从手机指向到眼镜观察的原位交互。方向性仍只是注意力的代理，下一节因而转向用户的指向、头向与注视如何贯穿对象发现、选择和反馈。

\section{注意力驱动的直观交互}

注意力驱动交互试图利用用户已经自然表达的位置、朝向、指向与注视，避免用户在物理对象和数字界面之间重复说明目标。在多对象环境中，一个完整过程至少包括候选对象发现、物理对象与数字身份绑定、动作确认和状态反馈。相关工作可据此分为三类：以视觉或无线信号发现对象，以空间输入选择对象，以及在虚实界面中维持对象状态。它们的差异不只在所用模态，更在于注意线索能够贯穿交互链条的哪一段。

第一类方法通过外部标记、外观或语义模型识别物体。二维码与人工视觉标签能够同时编码身份并提供几何定位，结果明确但要求标记处于视野内；自然特征法避免附加标记，却容易受弱纹理、相似外观、运动模糊和遮挡影响\supercite{wagnertracking}。de Freitas et al. 的 Snap-To-It 允许用户拍摄目标后在云端匹配，接近“拍到即选择”，但依赖预先准备的视觉数据库\supercite{snaptoit}。2024 年，Xu et al. 将几何、语义和视觉语言特征融合进三维场景表示，支持在 AR 中以自然语言搜索现实对象\supercite{spatial_ai_2024}；这类开放词汇方法扩展了可识别对象范围，却不能天然给出设备地址、访问权限和实时状态。

\begin{figure}[htbp]
  \centering
  \fbox{\parbox[c][4.0cm][c]{0.86\textwidth}{\centering
  \textbf{图片占位：现实对象发现方法谱系}\\[0.6em]
  建议内容：人工标记、自然特征、云端图像匹配与开放词汇三维理解，\\
  对比身份明确性、预部署需求、遮挡鲁棒性和设备地址获取能力。}}
  \caption{注意力驱动交互中的现实对象发现方法（占位图）}
  \label{fig:physical-object-discovery}
\end{figure}

第二类方法利用无线测位与空间传感缩小候选范围。Huo et al. 的 Scenariot 将超宽带定位与 SLAM 结合，把智能物体注册进 AR 场景\supercite{scenariot}；Zhang et al. 的 BLEselect 在智能眼镜上利用蓝牙到达角估计感知手势方向，把空间指向与无线身份关联\supercite{bleselect}；Kim et al. 的 IRIS 结合无线指环、手机图像和惯性数据，实现面向智能家居的指向控制\supercite{iris}。这些系统说明射频也能借助阵列或异构传感器获得方向，但通常要求设备预注册、环境建图或维护空间坐标与网络地址的映射，快速切换目标时仍可能出现绑定断点。

第三类工作关注虚拟表示与物理状态的持续一致。Zhu et al. 的 BISHARE 支持手机与头戴式 AR 之间的双向绑定，使手机既被头显感知又可控制虚拟内容\supercite{bishare}；Kaimoto et al. 将虚拟草图转化为实体机器的动作\supercite{sketchedreality}；Zhu et al. 则在 AR 界面与可动玩具之间建立实时映射\supercite{mecharspace}。这类系统将研究从“识别一个对象”推进到“让对象持续可控”，但其重点是已注册对象间的同步，首次接近陌生对象时仍需另行完成发现和授权。

\begin{figure}[htbp]
  \centering
  \fbox{\parbox[c][4.0cm][c]{0.86\textwidth}{\centering
  \textbf{图片占位：对象选择与虚实双向绑定}\\[0.6em]
  建议内容：无线测位/指向缩小候选对象，AR 界面绑定数字身份，\\
  再由无线会话同步物理状态与虚拟反馈。}}
  \caption{从空间对象选择到虚实状态同步的典型流程（占位图）}
  \label{fig:selection-and-bidirectional-binding}
\end{figure}

注意线索能够降低上述过程的显式开销。Pfeuffer et al. 的 gaze+pinch 以凝视定位、捏合确认，确立了 XR 中“快速指向+可靠确认”的常见分工\supercite{gaze_pinch}；其后续总结指出，凝视的速度优势必须与精度、舒适度、反馈和误触发共同权衡\supercite{gaze_pinch_challenges}。Schreiter et al. 对导航、探索与物体操作中的头眼协调进行分析，发现头向与视线的关系随活动而变化，因此不能把头向稳定地视作精确注视\supercite{human_gaze_head_2024}。Wilson et al. 进一步把眼动作为向情境 AI 表达注意的输入，但也强调注意信号需要结合任务上下文解释\supercite{gaze_contextual_ai}。这些近期结果共同表明，“方向更接近目标”并不自动等价于“用户要对目标执行动作”。

从这一视角看，RetroAR、EchoSight 并非孤立的两套光通信系统，而是对注意线索利用程度的两级推进。RetroAR 要求用户主动举起并指向手机，动作显著、意愿较强，代价是占用双手且需要对准\supercite{retroar}；EchoSight 让眼镜朝向自然跟随观察方向，并把身份返回和控制下行维持在同一局部光链路中，减少了设备列表选择和二维码扫描\supercite{echosight}。后者提高了自然度，却也削弱了“主动举机”所隐含的确认作用：用户可能只是阅读或扫视，系统不能仅凭眼镜朝向断定操作意图。

\begin{figure}[htbp]
  \centering
  \fbox{\parbox[c][4.2cm][c]{0.86\textwidth}{\centering
  \textbf{图片占位：注意力驱动的多对象交互链}\\[0.6em]
  建议内容：指向/头向/视线 $\rightarrow$ 候选目标 $\rightarrow$ 数字身份\\
  $\rightarrow$ 意图确认 $\rightarrow$ 会话与反馈，并标出信息断点。}}
  \caption{注意力线索在对象绑定与双向交互中的作用（占位图）}
  \label{fig:attention-interaction-chain}
\end{figure}

\medskip
\noindent\textbf{本节小结。}视觉、测位和虚实同步研究分别改善了对象发现、空间选择与状态保持，RetroAR/EchoSight 的改进是让同一方向线索贯穿“选中—识别—控制”。但 EchoSight 仍以机身朝向近似注视，且注视本身不等于操作；下一节因此讨论用眼动精确定位目标、再以生理证据判别意图的可能性。

\section{基于眼动、脑电信号的意图识别}

眼动是空间注意最直接的可观测线索之一。Jacob et al. 早期即使用视线指定界面元素，并指出眼动交互必须处理“看到即触发”的米达斯接触问题\supercite{jacobgaze}。与手持指向和头向相比，视线转移快、幅度小且不占用双手，适合从密集视野中快速给出候选目标；但它同时承担自然观察功能，系统不能把所有注视都解释为命令。

眼动首先受到测量误差和对象映射误差的限制。Feit et al. 对日常凝视输入的研究表明，准确度与精密度会随用户、屏幕区域和标定状态变化，并直接限制可可靠区分的目标尺寸与间距\supercite{everydaygaze}；Liu et al. 比较现实对象的 gaze-to-object 映射算法，也说明从二维注视点到三维物体的归属并非简单相交测试\supercite{gaze_to_object_mapping}。在 AR 眼镜中，设备滑移、头部运动、深度估计和遮挡还会持续改变眼动坐标、头部坐标与世界坐标之间的标定关系。因此，Omnisight 若以眼动驱动可调光束，首先需要验证的是动态条件下的闭环对齐误差和相邻目标可分性，而不是只报告离线眼动精度。

\begin{figure}[htbp]
  \centering
  \fbox{\parbox[c][4.0cm][c]{0.86\textwidth}{\centering
  \textbf{图片占位：从眼动测量到现实对象映射}\\[0.6em]
  建议内容：眼球坐标、头部坐标与世界坐标的转换链，\\
  标出标定误差、设备滑移、深度歧义和相邻目标混淆。}}
  \caption{眼动驱动现实目标定位中的坐标映射与误差来源（占位图）}
  \label{fig:gaze-to-object-error-chain}
\end{figure}

确认机制用于解决“看见”与“选择”的二义性。凝视停留简单但会牺牲速度并增加视觉负担，凝视加手势、按键或语音则将目标定位与命令确认分开\supercite{gaze_pinch,gaze_pinch_challenges}。这种显式确认可靠且易于理解，应当作为生理意图识别的强基线；只有当双手被占用、必须安静或频繁确认造成明显负担时，引入脑电才具有实际意义。换言之，眼动负责回答“看向哪里”，脑电至多提供“是否存在任务相关心理活动”的附加证据，两者不应被写成可直接读取完整意图。

脑电图（Electroencephalography, EEG）具有无创和时间分辨率高的特点，相关脑机接口大致分为诱发响应与自发活动两类。Farwell et al. 的经典 P300 拼写器利用用户关注低概率刺激时产生的事件相关电位确定目标\supercite{farwellp300}；Gao et al. 总结的 SSVEP 方法则依据注视周期刺激产生的稳态频率响应进行选择\supercite{visualauditorybci}。这两类范式目标明确、可形成较高信息传输率，但通常依赖闪烁刺激、重复试次或时间锁定平均，更适合离散界面选择，而不等价于自然环境中自发产生的操作意图。Abiri et al. 的综述也指出，可穿戴 EEG 仍受低信噪比、非平稳性、个体差异和伪迹影响\supercite{eegbcisurvey}。

\begin{figure}[htbp]
  \centering
  \fbox{\parbox[c][4.0cm][c]{0.86\textwidth}{\centering
  \textbf{图片占位：脑电意图识别范式比较}\\[0.6em]
  建议内容：P300、SSVEP 与自发 EEG 的刺激方式、时间窗口和输出，\\
  对比目标可解释性、响应速度、训练需求与自然交互适用性。}}
  \caption{诱发式与自发式 EEG 意图识别范式（占位图）}
  \label{fig:eeg-intention-paradigms}
\end{figure}

国内研究较早探索了眼动与 EEG 的互补。Wang et al. 在《航空学报》提出决策级融合方法，用眼动的空间与行为特征补充 EEG，对多类人机交互意图的识别优于单模态\supercite{wang_eeg_eye_2021}。Miao et al. 在 2024 年进一步以被动脑机接口建模五类隐式交互意图，说明 EEG 可提供浏览、搜索和处置任务中的区分信息，但其实验仍依赖实验室级 64 导设备与受控界面\supercite{miao_implicit_eeg_2024}。这些工作验证了“额外生理证据可能有益”，也暴露出从离线分类准确率到可穿戴实时交互之间的距离。

2024 年以来，多模态研究开始强调同步建模和自适应融合。Yi et al. 将 P300 与瞳孔收缩、视线响应输入双流多尺度网络，为意识障碍患者提供是/否通信；结果显示多模态系统可补充单一行为量表，但任务仍由强刺激和二选一答案严格约束\supercite{yi_hybrid_bci_2024}。Khanna et al. 在人机交接任务中直接比较 EEG、视线和手部运动，发现视线能够最早且最准确地预测交接意图，而融合主要提升较弱模态\supercite{khanna_handover_2025}。Zhu et al. 于 2026 年进一步以跨模态注意和动态集成融合 EEG、眼动与手势，实现六类交互意图分类\supercite{zhu_multimodal_intent_2026}。这些进展支持按场景动态选择模态，却也提醒系统不能预设 EEG 必然优于低成本行为信号。

对于本文拟开展的 Omnisight，上述文献给出了一条可检验而非预设成功的路线：先以眼动把 EchoSight 的“眼镜大致朝向”推进为“实测视线驱动光束”，完成现实目标与光学身份的精确对齐；再在观看但不操作、主动选择和执行控制等平衡任务中，比较眼动单模态、眼动+显式确认、EEG 单模态和眼动+EEG。评价指标除准确率外，还应包括误触发率、漏检率、决策时延、跨用户泛化、标定成本和佩戴负担。若 EEG 不能在这些端到端指标上稳定优于显式确认，则应保留眼动+低负担确认的退化路径，而不能仅凭离线分类结果宣称“所思即所控”。

\begin{figure}[htbp]
  \centering
  \fbox{\parbox[c][4.2cm][c]{0.86\textwidth}{\centering
  \textbf{图片占位：眼动与 EEG 协同的意图识别框架}\\[0.6em]
  建议内容：眼动给出空间目标，EEG 给出任务相关证据，\\
  置信融合后输出“观察/准备操作”，再决定执行或请求确认。}}
  \caption{眼动定位与脑电辅助判别的分级意图识别框架（占位图）}
  \label{fig:gaze-eeg-intention}
\end{figure}



